% Part: lambda-calculus
% Chapter: introduction
% Section: reduction

\documentclass[../../../include/open-logic-section]{subfiles}

\begin{document}

\olfileid{lam}{int}{red}
\olsection{લૅમ્બડા પદોનું લઘુકરણ}

લૅમ્બડા પદો સાથે શું કરી શકાય? તેમને સરળ બનાવી શકાય.
$M$ અને $N$ ગમે તે લૅમ્બડા પદો તથા~$x$ ગમે તે ચલ હોય, તો~$M$માં
$x$ની જગ્યાએ~$N$ મૂકવાથી મળતું પદ~$\Subst{M}{N}{x}$ વડે
દર્શાવી શકીએ. મૂક્યા પછી~$N$ના મુક્ત ચલો સાથે અથડામણ થાય એવા
$M$ના બંધ ચલોનું નામ પહેલેથી બદલી લેવું પડે. દાખલા તરીકે,
\[
\Subst{(\lambd[w][xxw])}{yyz}{x} = \lambd[w][(yyz)(yyz)w].
\]

\begin{digress}
સ્થાનાપન્ન માટેના બીજા સંકેતો $[N/x]M$, $[x/N]M$ અને
$M[x/N]$ પણ છે. ગૂંચવાડાથી સાવધાન!
\end{digress}

સહજ અર્થમાં $(\lambd[x][M])N$ અને~$\Subst{M}{N}{x}$નો અર્થ
સમાન છે. પહેલું પદ બીજા પદથી બદલવાની ક્રિયા
\emph{$\beta$-સંકોચન} કહેવાય છે. $(\lambd[x][M])N$ને
\emph{સંકોચ્ય પદ (રેડેક્સ)} અને~$\Subst{M}{N}{x}$ને તેનું
\emph{સંકોચન-ફળ} કહે છે. સામાન્ય રીતે, પદ~$P$ના કોઈ ઉપપદનું
$\beta$-સંકોચન કરીને~$P'$ મળે તો~$P$નું~$P'$માં
\emph{એક પગલામાં $\beta$-લઘુકરણ} થાય છે એમ કહીએ અને
$P \redone P'$ લખીએ. $P$માંથી એક-પગલાના અમુક લઘુકરણો દ્વારા
(શૂન્ય પણ હોઈ શકે) $P'$ મળે તો~$P$નું~$P'$માં
\emph{$\beta$-લઘુકરણ} થાય છે; સંકેતમાં~$P \red P'$.
જે પદનું આગળ~$\beta$-લઘુકરણ થઈ ન શકે તેને
\emph{$\beta$-અલઘુકરણીય} અથવા \emph{$\beta$-સામાન્ય}
કહે છે. સંદર્ભ સ્પષ્ટ હોય ત્યારે ``$\beta$-લઘુકરણ''ને બદલે
માત્ર ``લઘુકરણ'' વગેરે કહીશું.

કેટલાંક ઉદાહરણો જોઈએ.
\begin{enumerate}
\item આપણને મળે છે:
\begin{align*}
(\lambd[x][xxy]) \lambd[z][z] & \redone (\lambd[z][z])(\lambd[z][z]) y \\
& \redone (\lambd[z][z]) y \\
& \redone y.
\end{align*}
\item પદને ``સરળ'' કરતાં તે વધુ જટિલ પણ બની શકે છે:
\begin{align*}
(\lambd[x][xxy])(\lambd[x][xxy]) & \redone (\lambd[x][xxy])(\lambd[x][xxy])y \\
& \redone (\lambd[x][xxy])(\lambd[x][xxy])yy \\
& \redone \dots
\end{align*}
\item લઘુકરણ પછી પદ યથાવત્ પણ રહી શકે છે:
\[
(\lambd[x][xx])(\lambd[x][xx]) \redone (\lambd[x][xx])(\lambd[x][xx]).
\]
\item વળી, કેટલાક પદોનું એક કરતાં વધુ રીતે લઘુકરણ થઈ શકે છે.
  દાખલા તરીકે,
\[
(\lambd[x][(\lambd[y][yx]) z)] v \redone (\lambd[y][yv]) z
\]
બહારના અનુપ્રયોગનું સંકોચન કરીને મળે છે; અને
\[
(\lambd[x][(\lambd[y][yx]) z)] v \redone (\lambd[x][zx]) v
\]
અંદરના અનુપ્રયોગનું સંકોચન કરીને મળે છે. જોકે આ કિસ્સામાં
બંને પદો આગળ જઈને એક જ પદ~$zv$માં લઘુકૃત થાય છે.
\end{enumerate}

છેલ્લા ઉદાહરણમાં મળેલું સમાન અંતિમ પરિણામ સંયોગ નથી. તે
લૅમ્બડા કલનનો ``ચર્ચ--રોસર ગુણધર્મ'' નામનો ઊંડો અને મહત્વનો
ગુણધર્મ બતાવે છે.

\end{document}
